
faisant intervenir les anneaux de classes de cycles (mod. équivalence rationnelle) @@M0@@ et @@M1@@, commutativité établie dans SGA 5 IV.

\textbf{6.3.} Notons que dans chacun des cas envisagés plus haut où on arrive à définir un homomorphisme de Gysin, on obtient en fait une construction pour la cohomologie à coefficients dans @@M2@@, et non seulement dans @@M3@@. Cela permet, lorsque @@M4@@ est une \uline{immersion}, de poser la question de la validité de la formule de Riemann--Roch ``sans dénominateurs'', comme il est expliqué au n° 3. Lorsque @@M5@@ et @@M6@@ sont tous deux lisses sur un corps @@M7@@, la commutativité de (6.4) nous montre qu'il suffirait de prouver la formule sans dénominateurs dans l'anneau de Chow @@M8@@ cf. (3.3).

\textbf{6.4.} Lorsqu'on ne fait pas d'hypothèse de régularité sur @@M9@@, et même si @@M10@@ et @@M11@@ sont de type fini sur le corps @@M12@@, on ne voit plus de définition plausible d'un homomorphisme de Gysin cohomologique (6.2), même pour une immersion régulière de codimension @@M13@@. Peut-on en trouver un néanmoins avec des propriétés raisonnables, en particulier rendant commutatif le carré (6.3)? La question est liée évidemment à celle de définir, pour un @@M14@@ d'intersection complète de codimension @@M15@@ dans @@M16@@, ``sa classe de cohomologie''
@@M17@@
de façon à transformer intersections (sans composantes excédentaires) en cup-produits. Cette dernière question se résout affirmativement tout au moins modulo torsion, i.e. dans @@M18@@, en définissant
@@M19@@
Mais j'ignore si @@M20@@ est divisible par @@M21@@ dans @@M22@@. En particulier, si @@M23@@ est un schéma projectif de dimension @@M24@@ sur le corps des complexes, disons, et @@M25@@ réduit à un point (singulier) de @@M26@@, @@M27@@ est-il divisible par @@M28@@ dans @@M29@@, et en particulier la classe de Chern est-elle nulle dans @@M30@@? Signalons également dans la même voie la question suivante: soit @@M31@@ un schéma projectif (singulier) de dimension @@M32@@, @@M33@@ un sous-schéma fermé de dimension @@M34@@ qui soit globalement une intersection de trois diviseurs positifs @@M35@@, considérons le produit de leurs classes de Chern mod. @@M36@@:
@@M37@@
